% Convolutional Neural Networks — see issue #25.
\subsection{Convolutional Neural Networks}
\label{subsec:cnns}

A \emph{convolutional neural network} (CNN) is a neural network that uses convolutions in at least some of its layers. A \emph{convolution} is an operation in which a small grid of weights, called a \emph{kernel} or \emph{filter}, is moved across the input and, at each position, multiplied element by element with the values it covers and summed. A \emph{convolutional layer} is a layer that applies one or more kernels to its input; each kernel is small, typically $3 \times 3$ in 2D or $3 \times 3 \times 3$ in 3D, and its weights are learned during training. These characteristics make CNNs well suited to data with a grid structure, such as images \parencite{prince2023understanding}.

For a 2D image $X$ and a kernel $K$, the output of a convolutional layer at position $(i, j)$ is given by \cref{eq:convolution}, where $b$ is a bias term and $\sigma$ is a non-linear activation function, for example ReLU, which replaces negative values with zero:
\begin{equation}
  F(i,j) = \sigma\!\left( \sum_{m}\sum_{n} K(m,n)\, X(i+m,\, j+n) + b \right)
  \label{eq:convolution}
\end{equation}
The output obtained by applying one kernel to the whole input is called a \emph{feature map}, and it shows where in the input the pattern encoded by that kernel appears. A layer applies many kernels in parallel, so its output is a stack of feature maps that becomes the input of the next layer.

To reduce the size of the feature maps, convolutional layers are usually followed by \emph{pooling}, an operation that replaces each small neighborhood of a map with a single value, either its maximum (max pooling) or its average (average pooling), which also makes the representation less sensitive to small shifts of a pattern. Both the kernel and the pooling window are controlled by two parameters:
\begin{itemize}
  \item \textbf{Stride:} the number of positions the window advances at each step; a stride larger than 1 also shrinks the output.
  \item \textbf{Padding:} values, usually zeros, added around the border of the input so that the output keeps the same size as the input.
\end{itemize}

Because each layer takes the output of the previous one as its input, a unit deep in the network depends on a growing region of the original input, called its \emph{receptive field}. As a result, early layers detect simple patterns such as edges and textures, and later layers detect larger structures built from them.

The same operations extend to three dimensions. In a 2D CNN the kernel moves along the height and width of an image, whereas in a 3D CNN it also moves along the depth, so it processes a whole volume and preserves the relation between neighboring slices, at the cost of many more parameters and memory.

Beyond these basic operations, CNNs differ in how many layers they stack and how those layers are connected. Three architectures became standard references \parencite{prince2023understanding}, summarized in \cref{tab:cnn_architectures}.

% Reference CNN architectures — used in sections/theoretical_framework/cnns.tex.
\begin{table}[t]
  \centering
  \caption{Reference CNN architectures.}
  \label{tab:cnn_architectures}
  \small
  \begin{tabular}{@{}p{0.2\columnwidth}p{0.22\columnwidth}p{0.5\columnwidth}@{}}
    \toprule
    Architecture & Depth & Key idea \\
    \midrule
    VGG      & 16 or 19 layers  & Stacks of small $3 \times 3$ kernels; simple and deep. \\
    ResNet   & 18 to 152 layers & Skip connections, shortcuts that add a block's input to its output, make very deep networks trainable. \\
    DenseNet & 121 to 201 layers & Each layer receives the feature maps of all previous layers, which reuses features and needs fewer parameters. \\
    \bottomrule
  \end{tabular}
\end{table}

